% Metodologia do artigo (template SBC) — Lab03S02, Issue #31. RASCUNHO baseado só em artefatos do repositório.
% Cada afirmação factual aponta a fonte em comentário (%). Itens que dependem do professor aparecem como
% \pend{...}; decisões provisórias do grupo como \prov{...}. Evidências e lacunas: docs/metodologia_evidencias.md.
\newcommand{\pend}[1]{\textbf{[PENDENTE DO PROFESSOR: #1]}}
\newcommand{\prov}[1]{\textbf{[PROVISÓRIO: #1]}}

\section{Metodologia}

\subsection{Questões de pesquisa e desenho do estudo}
% E§5; fonte_vdd.md
O estudo é observacional e mineratório: calcula as métricas DORA a partir de dados públicos do GitHub
(\textit{releases}, \textit{commits} e execuções de GitHub Actions) de repositórios open-source. As
questões de pesquisa seguem o enunciado: RQ01 frequência de deploy; RQ02 \textit{lead time} (variantes
por \textit{release} e por \textit{commit}); RQ03 \textit{change failure rate} (variante (a), proxy de CI,
e (b), proxy de entrega); RQ04 tempo de recuperação; RQ05 correlação entre frequência e CFR; RQ06
características associadas ao desempenho; RQ07 sensibilidade da classificação DORA à definição operacional.
Como o GitHub não registra deploys nem falhas em produção, todas as medidas são \emph{proxies}.
Esta seção descreve a coleta e a validação manual; as análises estatísticas das RQ01--RQ07 pertencem à
Sprint~3 e \textbf{ainda não foram executadas}.

\subsection{Seleção dos repositórios e funil}
% outputs/funil.md, config.yaml (selecao:)
Os candidatos vieram da Search API, fatiada por faixas de estrelas a partir de 1000 (\texttt{stars:>=1000},
sem forks e sem repositórios arquivados). \prov{o enunciado não define ``popular''; o piso de 1000 estrelas,
a exclusão de forks/arquivados e as faixas de fatiamento são decisão do grupo (lacuna L03)}.
Critérios de inclusão (E§3): uso de GitHub Actions, ao menos 5 \textit{releases} e ao menos 50
\textit{workflow runs} válidos no \textit{default branch} dentro da janela. \prov{a leitura de ``run válido''
(\texttt{event=push}, criado na janela, \texttt{conclusion} em \{success, failure, timed\_out,
startup\_failure\}) e a de ``release contada'' (não rascunho, não pré-release, publicada na janela) são
interpretações do grupo (L04, L05)}.

\begin{table}[ht]
\centering
\caption{Funil de seleção (\texttt{outputs/funil.csv}).}
\begin{tabular}{lrrr}
\hline
Etapa & Entram & Saem & Restam \\ \hline
Candidatos avaliados (sem duplicatas) & 2677 & --- & 2677 \\
Possui Actions & 2677 & 1219 & 1458 \\
Releases mínimas (5) & 1458 & 1038 & 420 \\
Runs válidos mínimos (50) & 420 & 120 & 300 \\
Amostra final & 300 & --- & 300 \\ \hline
\end{tabular}
\end{table}

A Search API devolveu 60\,143 resultados brutos (60\,142 após deduplicação), dos quais 2677 foram
avaliados em ordem determinística por \texttt{full\_name}; a avaliação parou ao atingir 300 repositórios,
e 57\,465 candidatos \emph{não foram avaliados}. Assim, a amostra não é aleatória sobre a população de
candidatos: é o primeiro conjunto de 300 que passa nos filtros nessa ordem (lacuna L17, corte não definido
pelo enunciado). O relatório marca todas as fatias como não saturadas, mas a fatia \texttt{stars:1282..1312} tem
exatamente 1000 resultados (o teto da API); \pend{confirmar se esse caso foi truncado}.

\subsection{Janela temporal}
% config.yaml (L01); data/amostra_ouro/sorteio.json
A janela de 12 meses usada nos artefatos atuais é 2025-10-02 a 2026-10-02 (UTC). \pend{as datas oficiais
da janela não constam do material fornecido (lacuna L01); o \texttt{config.yaml} as declara
\emph{provisórias}. Se mudarem, seleção, coleta e sorteio da amostra-ouro precisam ser refeitos}.

\subsection{Definições operacionais das métricas}
% E§3, E§5; metricas/*.py; README
Valem as definições comuns do enunciado: deploy é uma \textit{release} publicada (\texttt{draft=false},
sem pré-release na definição principal); a data de um \textit{commit} é \texttt{commit.author.date};
considera-se apenas o \textit{default branch} para execuções de CI, com \texttt{event=push}; a
classificação de falha/sucesso segue a tabela de \texttt{conclusion} do enunciado, ignorando
\texttt{cancelled}, \texttt{skipped}, \texttt{neutral}, \texttt{action\_required}, \texttt{stale} e execuções
em andamento.
\begin{itemize}
  \item \textbf{Deployment frequency:} \textit{releases} publicadas na janela dividido pelo número de
  semanas da janela (RQ01).
  \item \textbf{Lead time (a)/(b):} por \textit{release}, \texttt{publicação $-$ commit mais antigo}; por
  \textit{commit}, \texttt{publicação $-$ data do commit}; mediana por repositório, em horas. Commits
  vêm de \texttt{compare/\{anterior\}...\{R\}}; a primeira \textit{release} da história é ignorada, e
  \textit{releases} com \texttt{compare} 404 também. \prov{releases sem commits novos e commits com data
  posterior à release são excluídos (L07)}.
  \item \textbf{CFR (a):} falhas $\div$ (falhas $+$ sucessos) das execuções de CI.
  \item \textbf{CFR (b):} uma \textit{release} R falha se qualquer \textit{release} de deploy em
  $(R, R{+}7\text{d}]$ for corretiva; \textit{releases} nos últimos 7 dias da janela ficam fora do
  denominador (censura); com menos de 2 \textit{releases} na janela o valor é indefinido.
  \prov{``qualquer release no intervalo'', e não só a próxima, é decisão do grupo (L08)}.
  \item \textbf{Tempo de recuperação:} por workflow, do início da primeira falha após um sucesso até o fim
  (\texttt{updated\_at}) do próximo sucesso; mediana por repositório; episódios sem fim na janela são
  censurados e reportados em proporção.
  \item \textbf{Classificação DORA:} cortes fixos adaptados pelo enunciado (Elite/High/Medium/Low) e
  mediana arredondada para baixo das quatro notas.
\end{itemize}
O CFR (b) \textbf{não} está no arquivo de métricas atual (\texttt{outputs/resultados.csv}): a coluna
\texttt{cfr\_b} não existe nele, porque a heurística não foi validada (Seção~\ref{sec:heuristica}).

\subsection{Procedimento de coleta, cache e limitações da API}
% README; pipeline/http_client.py; outputs/execucao.log, funil.md
A coleta é feita por script próprio, via API REST, sem bibliotecas de acesso ao GitHub. Cada resposta é
armazenada em disco e a execução é retomável: o último \textit{run} registrado em
\texttt{outputs/execucao.log} processou 300 repositórios sem erros, com 12\,191 acertos e 72\,975 falhas de
cache (execução iniciada em 2026-10-11). O cliente espera o \textit{rate limit} renovar (até 10 esperas por
requisição) e repete erros 5xx com \textit{backoff} exponencial (1, 2, 4, 8~s, até 5 tentativas).
Limitações conhecidas: (i) a Search API devolve no máximo 1000 resultados por consulta, por isso a busca é
fatiada; (ii) execuções de CI são coletadas por mês, e nenhum dia atingiu o teto de 1000 resultados
(\texttt{n\_runs\_dias\_saturados} $=0$ em todos os repositórios de \texttt{resultados.csv}); (iii) o
\texttt{compare} é paginado; (iv) \textit{releases} com \texttt{compare} divergente (backports) mantêm
os commits desde o \textit{merge-base} e podem inflar o \textit{lead time}; (v) o cache não é versionado,
portanto a coleta só é reproduzível por quem o possui ou refaz as chamadas, e \emph{repositórios continuam
recebendo commits}, de modo que uma nova coleta pode divergir.

\subsection{Rotulagem manual, sugestões assistidas e consenso}
% data/amostra_ouro/{sorteio.json,LEIAME.md,PROTOCOLO_CONSENSO.md,consenso.csv}; rascunho/revisao_*.csv
Foram sorteados 60 repositórios da lista de entrada \texttt{resultados.csv} (semente 42; \texttt{sha256}
registrado em \texttt{sorteio.json}) e 5 \textit{releases} por repositório (300 \textit{releases}, as
mesmas para os três avaliadores). Cada avaliador preencheu, em arquivo próprio, (i) o tipo do projeto,
(ii) se as \textit{releases} são entregas reais e (iii) se cada \textit{release} é corretiva.

\textbf{Uso de sugestões.} A rotulagem usou sugestões: \texttt{scripts/rotular.py} apresenta notas e
commits de cada item junto de uma \emph{sugestão}, igual para todos os avaliadores, que o avaliador aceita
ou substitui. Existem 292 sugestões para os 300 itens de \texttt{corretiva} (71 ``sim'', 221 ``não''); os
8 restantes ficaram sem sugestão.

\textbf{Rodada anterior (rótulos que não são os commitados).} Numa primeira rodada, feita com versões dos
arquivos de rótulos que não foram as commitadas, os arquivos de revisão de \textbf{dois} avaliadores
(guardados fora do Git) coincidiam com a sugestão em \textbf{292/292} itens, esses dois avaliadores
coincidiam entre si em 293 dos 300 itens (97,7\%) e o terceiro coincidia com a sugestão em 181 dos 292
(62,0\%).

\textbf{Arquivos commitados (usados nos resultados abaixo).} Os rótulos oficiais são os três
\texttt{labels\_<avaliador>.csv} commitados nas Issues \#23, \#24 e \#25. Neles, a concordância em
\texttt{corretiva} é de 82,0\% (246/300) entre ALuiz19 e marcosffp, 63,3\% (190/300) entre ALuiz19 e
GuilhermeVieira05 e 56,0\% (168/300) entre GuilhermeVieira05 e marcosffp. A aderência desses arquivos à
sugestão não pode ser recalculada, porque \texttt{sugestoes.csv} e os arquivos de revisão estão fora do Git.

Não há, nos artefatos, registro de como as sugestões foram geradas (\texttt{sugestoes.csv} traz só rótulo
e motivo textual). Os fatos acima não permitem concluir que os rótulos estejam errados: uma sugestão
pode estar correta. Eles apenas \emph{não demonstram} três julgamentos independentes, e o critério de
independência do enunciado (``sem consultar os demais'') não é verificável nesses casos.

\textbf{Repositório removido.} O repositório \texttt{1bananachicken/MaaNTE} (5 \textit{releases} da
amostra) foi removido do GitHub depois do sorteio: a API devolve 404 para o repositório, as
\textit{releases}, os \texttt{compare} e a conta do dono. Ele foi mantido na amostra por decisão do grupo.
Dois dos três avaliadores rotularam essas \textit{releases} sem acesso ao conteúdo (registrado na coluna
\texttt{observacoes}).
\pend{se a rotulagem assistida é metodologicamente aceitável para a validação, ou se a amostra deve ser
rotulada de novo sem sugestões, é decisão do professor}.

O protocolo de consenso (\texttt{PROTOCOLO\_CONSENSO.md}) é: unanimidade; maioria simples em 2 a 1;
reunião para 1-1-1. O protocolo prevê o kappa antes do consenso. Resultado em \texttt{consenso.csv}:
\texttt{corretiva}, 152 por unanimidade e 148 por maioria (85 corretivas, 215 não); tipo do projeto, 44 e
16; entrega real, 55 e 5; nenhum caso \texttt{reuniao} nem \texttt{pendente}.

\subsection{Concordância entre avaliadores}
% data/amostra_ouro/concordancia.csv; recalculado a partir de labels_*.csv com statsmodels
\begin{table}[ht]
\centering
\caption{Kappa de Fleiss (3 avaliadores). Interpretação de Landis \& Koch.}
\begin{tabular}{lrrl}
\hline
Dimensão & Itens & $\kappa$ & Interpretação \\ \hline
Tipo do projeto & 60 & 0,743 & substancial \\
Releases são entregas reais & 60 & 0,352 & razoável \\
Release corretiva & 300 & 0,273 & razoável \\ \hline
\end{tabular}
\end{table}
O kappa baixo da dimensão \texttt{corretiva} indica que o próprio critério é ambíguo, o que é um achado e
limita qualquer validação contra esse consenso. A concordância simples foi 0,507 em \texttt{corretiva}.
Em \texttt{entrega real}, a concordância simples é alta (0,917), mas quase todos os rótulos são ``sim'',
o que derruba o kappa (efeito de prevalência).
\textbf{Limitação na leitura do kappa.} O kappa de Fleiss pressupõe avaliações independentes, o que os
artefatos não demonstram (ver acima). Com a maioria 2 a 1, o par de avaliadores mais próximo (82,0\% de
concordância em \texttt{corretiva}) pesa mais no consenso do que o terceiro. Para a heurística, o consenso
usado como referência de precisão, \textit{recall} e F1 pode herdar o viés das sugestões; o sentido e a
magnitude desse efeito não foram medidos. Os valores foram calculados por \texttt{scripts.agreement} a
partir dos três arquivos commitados (\texttt{concordancia.csv}).
Outras medidas de avaliação usadas: precisão, \textit{recall} e F1 (\texttt{precision\_recall\_fscore\_support},
positivo = corretiva). Nenhum intervalo de confiança foi calculado. \pend{se o professor exige
intervalos de confiança ou outro teste para essas medidas, não há procedimento definido nos documentos}.

\subsection{Heurística de release corretiva: versões e desempenho}\label{sec:heuristica}
% data/amostra_ouro/avaliacao_heuristica.csv; docs/refino_heuristica.md; git log
A heurística v1 reproduz o exemplo do professor (``só muda o \textit{patch} e contém \texttt{fix}'') e foi
definida antes de ver os dados. As versões v2--v7 combinam os mesmos dois sinais (SemVer e palavra-chave,
estrita ou ampla). As versões v8 e v9 olham o \emph{conteúdo} (commits de correção e de feature; na v9,
também as \textit{release notes}) e foram desenhadas \textbf{depois} de analisar os erros da v1--v7 na
\emph{mesma} amostra em que são avaliadas. Além disso, v8 e v9 entraram no repositório (PR \#65, 14/10,
21h19) \emph{antes} de os rótulos de dois avaliadores serem commitados (PRs \#67 e \#68, 21h46 e 21h59):
foram desenhadas sobre a rodada anterior de rótulos.

Cada versão é medida de duas formas: em todas as \textit{releases} em que ela é aplicável (300; a v9,
que exige as notas, em 295, porque as 5 \textit{releases} do repositório removido não têm notas) e no
\emph{subconjunto comum} às nove versões (295 \textit{releases}), usado para compará-las.

\begin{table}[ht]
\centering
\caption{Versões testadas contra o consenso (300 \textit{releases}; 85 corretivas). Precisão, \textit{recall}
e F1 em todas as \textit{releases} aplicáveis ($n$); F1 comum no subconjunto de 295. Fonte:
\texttt{avaliacao\_heuristica.csv}.}
\begin{tabular}{llrrrrrl}
\hline
Versão & Regra & $n$ & Precisão & Recall & F1 & F1 comum & Desenho \\ \hline
v1 & patch \textbf{e} palavra estrita & 300 & 0,423 & 0,741 & 0,539 & 0,546 & antes dos dados \\
v2 & patch \textbf{ou} palavra estrita & 300 & 0,329 & 0,988 & 0,494 & 0,485 & antes da análise de erros\textsuperscript{*} \\
v3 & só palavra estrita & 300 & 0,313 & 0,835 & 0,455 & 0,460 & idem\textsuperscript{*} \\
v4 & só patch (não comparável: palavra) & 300 & 0,429 & 0,894 & 0,580 & 0,570 & idem\textsuperscript{*} \\
v5 & patch \textbf{e} palavra ampla & 300 & 0,432 & 0,788 & 0,558 & 0,565 & idem\textsuperscript{*} \\
v6 & patch \textbf{ou} palavra ampla & 300 & 0,328 & 0,988 & 0,493 & 0,484 & idem\textsuperscript{*} \\
v7 & só palavra ampla & 300 & 0,321 & 0,882 & 0,470 & 0,475 & idem\textsuperscript{*} \\
v8 & correção e nenhuma feature & 300 & 0,625 & 0,412 & 0,497 & 0,507 & depois da análise de erros \\
v9 & v8 \textbf{ou} notas com correção e sem novidade & 295 & 0,627 & 0,573 & 0,599 & 0,599 & depois da análise de erros \\ \hline
\end{tabular}
\par{\footnotesize \textsuperscript{*} v2--v7 entraram no mesmo \textit{commit} do script de avaliação
(\texttt{360adc7}), antes dos arquivos de resultado; a ordem exata entre escrevê-las e a primeira rodada
não é verificável pelo histórico.}
\end{table}

\textbf{Status: heurística não validada.} O critério do enunciado é F1 $\geq$ 0,70, e o melhor F1 observado
foi 0,599 (v9, $n$ = 295); entre as versões aplicáveis a todo o \textit{dataset}, o melhor foi 0,580 (v4).
Conforme o enunciado, as versões e o F1 de cada uma são documentados aqui. A coluna
\texttt{escolhida} de \texttt{avaliacao\_heuristica.csv} é apenas a regra interna do script (maior F1 comum
quando ninguém atinge 0,70) e \emph{não} indica validação. \prov{o grupo adotou v8 como estimativa exploratória do
CFR (b), v4 como sensibilidade e v9 fora do dataset final (\texttt{docs/refino\_heuristica.md})}; com os
rótulos commitados, porém, a v8 (F1 0,497) ficou abaixo da v4 (0,580) e da v5 (0,558), e essa estratégia
precisa ser revista pelo grupo. A v9
depende das \textit{release notes} (\texttt{body}), que a coleta atual não armazena; por isso ela só foi
aplicada à amostra-ouro, onde as notas foram buscadas à parte. \pend{o enunciado e a Fonte da Verdade
não dizem como proceder se nenhuma versão chegar a F1 $\geq$ 0,70, nem se o CFR (b) exploratório é
aceitável; também não definem amostra independente para validação final}.

\subsection{Dados ausentes, censura e ameaças à validade}
% outputs/resultados.csv (contagem direta); E§3, E§5, E§11
\textbf{Dados ausentes e censura.} Casos censurados não são descartados: dos 9062 episódios de falha
em \texttt{resultados.csv}, 176 foram censurados; 6 \textit{releases} foram ignoradas no \textit{lead
time} por 404 no \texttt{compare}; 23 repositórios têm tempo de recuperação indefinido (sem episódios
concluídos). \prov{a mediana de recuperação usa só episódios concluídos e a proporção de censurados é
reportada à parte (L09)}. No CFR (b), \textit{releases} dos últimos 7 dias da janela ficam fora do
denominador. Em \texttt{resultados.csv}, 381 commits (32 repositórios) foram excluídos do \textit{lead time} por data posterior à release e 1685 \textit{releases} (51 repositórios) não têm commits novos (\texttt{python -m scripts.data\_quality}). Cinco \textit{releases} da amostra-ouro têm mensagens incompletas (\texttt{compare} 404 ou
truncado), o que pode ocultar a palavra-chave.

\textbf{Ameaças à validade} (resumo; a seção completa é entregue com o artigo final):
\begin{itemize}
  \item \emph{Construto:} \textit{release} $\neq$ deploy; falha de CI $\neq$ falha em produção;
  \textit{commit} com \textit{rebase}/\textit{squash} distorce a data; \textit{releases} não são
  filtradas pelo \textit{default branch}; a heurística corretiva não é validada.
  \item \emph{Interna:} v8/v9 foram desenhadas sobre a amostra em que são avaliadas, antes de os rótulos
  oficiais de dois avaliadores serem commitados, e a escolha é feita entre nove versões, o que torna o F1
  otimista; a rotulagem assistida reduz a independência dos avaliadores; um repositório da amostra foi
  removido e rotulado sem acesso ao conteúdo por dois avaliadores.
  \item \emph{Externa:} amostra dos primeiros 300 repositórios que passam nos filtros em ordem
  alfabética, com piso de 1000 estrelas; não representa projetos corporativos nem a população de
  candidatos.
  \item \emph{Conclusão:} kappa razoável (0,273) em \texttt{corretiva}; a amostra-ouro tem 60 repositórios e
  5 \textit{releases} por repositório; \textbf{nenhuma análise estatística das RQs foi feita até aqui}.
\end{itemize}
